\specialchapter{HISTORICAL NOTE}

(Chapters I to V)

The notions of eigenvalue and Fourier series were already well known in 1830, permeating many fields of analysis throughout the nineteenth century (cf.~ÉHM, pp.~114, 260, 275). But spectral theory and harmonic analysis, in the sense in which we understand them in this book, did not begin to take their present form until about a century later --- at the same time as their junction with the study of group representations around 1930, and then with that of normed algebras around 1940.

We confine ourselves here to sketching the main lines of their evolution in the period from 1906—when Hilbert introduced the notion of the “continuous spectrum” in his work on integral equations—to the years 1945–1950, when the theories presented here essentially acquired the form they have retained to the present day.

